\begin{definition}
  A matrix $Q \in M_n(\mathbb{R})$ is orthogonal if $Q^T Q = \mathbb{I}_n$
\end{definition}

\begin{remark}
    $Q$ is orthogonal if and only if the columns of $Q$ are orthonormal with respect to the dot product.
\end{remark}

\begin{theorem} (QR Decomposition)
  
  Let $A \in M_{n}(\mathbb{R})$ be invertible. Then there exists orthogonal $Q$ and upper-triangular $R$ such that
  \[A = QR.\]
\end{theorem}
\begin{proof}
    Let $A$ have columns $v_1,...,v_n$. Then as $A$ is invertible these are a basis of $\mathbb{R}^n$. Applying the Gram-Schmidt process gives and orthonormal basis $u_1,...,u_n$. Then
    \begin{align*}
      v_1 &= \norm{w_1} u_1 \\
      v_2 &= \norm{w_2} u_2 + \langle u_1, v_2 \rangle u_1 \\
      ..&. \\
      v_k &= \norm{w_k}v_k + \sum_{ j = 1}^{k} \langle u_j, u_k \rangle u_j
    .\end{align*}
    Then $u_1,...,u_n$ define the columns of $Q$ and let 
    \[R_{ij} = \left\{\begin{array}{lr}
    \langle u_i, v_j \rangle & i < j\\
    \norm{w_i} & i = j\\
    0 & i > j\\
    \end{array} \right..\]
    Then $A = QR$.
\end{proof}

\begin{remark}
  To compute $Q$ and $R$ in practise it is simplest to first find $Q$, then at $Q^T = Q^{-1}$, we have 
  \[R = Q^T A.\]
\end{remark}

\subsection{Orthogonal Complements}

Let $V$ be an inner product space and let $U \leq V$. Recall $u, v \in V$ are orthogonal, $u \perp v$, if $\langle u, v \rangle = 0$.

\begin{definition}
  The {\bf orthogonal complement} $U^\perp$ of $u$ in $V$ is
  \[U^\perp = \{ v \in V \ | \ \langle u, v \rangle = 0 \ \text{for} \ u \in U \}.\]
\end{definition}

\begin{prop} $ $

  \begin{itemize}
    \item $U^\perp$ is a subspace of $V$
    \item $U \cap U^\perp = \{0\}$
    \item $U \subseteq (U^\perp)^\perp$
  \end{itemize}
\end{prop}
\begin{proof} $ $

   \begin{itemize}
     \item From the linearity in the second argument
     \item Let $u \in U$ and $u \in U^\perp$. Then $\langle u, u \rangle = 0$ which implies $u=0$ by the positive definite axiom.
     \item If $u \in U, w \in U^\perp$ then
       \[\langle w, u \rangle = \overline{\langle u, w \rangle} = 0\]
       so $u \in (U^\perp)^\perp$.
   \end{itemize} 
\end{proof}

\begin{theorem}
    If $U \leq v$ is a finite dimensional subspace then $V = U + U^\perp$.
\end{theorem}
\begin{proof}
    It's sufficient to show for any $v \in V$ there exists $v_1 \in U$ and $v_2 \in U^\perp$ such that $v = v_1 + v_2$.
    
    \noindent Let $u_1,...,u_k$ be an orthonormal basis of $U$. Given $v \in V$ set
    \begin{align*}
      v_1 &= \sum_{ i = 1}^{k} \langle u_i, v \rangle u_i \\
      v_2 &= v - v_1
    .\end{align*}
    Then $v_1 \in U$ and to show $v_2 \in U^\perp$:
    \begin{align*}
      \langle u_j, v_2 \rangle &= \langle u_j, v \rangle - \langle u_j, v_1 \rangle \\
       &= \langle u_j, v \rangle - \sum_{ i = 1}^{k} \langle u_i, v \rangle \langle u_j, u_j \rangle \\
       &= \langle u_j, v \rangle - \langle u_j, v \rangle = 0
    .\end{align*}
\end{proof}

\begin{corollary}
    Let $V$ be a finite dimensional inner product space.

    \begin{itemize}
      \item $\dim U^\perp = \dim V - \dim U$ 
      \item $U = (U^\perp)^\perp$
    \end{itemize}
\end{corollary}
\begin{proof} $ $
  \begin{itemize}
    \item As $V = U + U^\perp$, we have
      \begin{align*}
        \dim V &= \dim U + \dim U^\perp - \dim (U \cap U^\perp) \\
               &= \dim U + \dim U^\perp - 0
      .\end{align*}
    \item $U \leq (U^\perp)^\perp$. As they have the same rank they are equal.
  \end{itemize}
\end{proof}

